\section{Conclusion and future work}

The suite walks a kernel ladder from a memory bound SAXPY to a compute bound
register blocked GEMM, places each rung on a roofline against both the
theoretical and the measured ceilings for this RTX 5070, and ties each position
back to the Nsight Compute counter that explains it. The value is not any single
fast kernel but the counter backed account of why each kernel lands where it
does, and the honest gap between the vendor spec sheet and what the card actually
delivers.

Several extensions would sharpen the picture. A WMMA tensor core GEMM on its own
roofline panel would add a second, higher compute ceiling and show how far the
FP32 CUDA core story is from the tensor core story on the same silicon. A cache
aware roofline with a second sloped ceiling for L2 bandwidth would separate
kernels that are DRAM bound from kernels that are L2 bound, which the single
DRAM ceiling here cannot distinguish. Sweeping the power and clock limits and
replotting would turn the NVML trace from a diagnostic into a first class axis,
showing how the ceilings themselves move as the card is throttled. Each of these
is additive: the harness, the analysis package, and the report pipeline are built
to absorb a new kernel or a new ceiling and regenerate every figure with one
command.
